% Neural Surrogate Modeling for Snake Robot Locomotion
% Research Report
% Phase 6: Write Research Report in LaTeX

\documentclass[11pt]{article}

% ---------------------------------------------------------------------------
% Font and encoding
% ---------------------------------------------------------------------------
\usepackage{palatino}
\usepackage[T1]{fontenc}
\usepackage[utf8]{inputenc}

% ---------------------------------------------------------------------------
% Page layout
% ---------------------------------------------------------------------------
\usepackage[top=1in, bottom=1in, left=1.25in, right=1.25in]{geometry}
\usepackage{setspace}
\setstretch{1.1}
\setlength{\parindent}{1em}
\setlength{\parskip}{0.4em}

% ---------------------------------------------------------------------------
% Math
% ---------------------------------------------------------------------------
\usepackage{amsmath}
\usepackage{amssymb}

% ---------------------------------------------------------------------------
% Figures and tables
% ---------------------------------------------------------------------------
\usepackage{graphicx}
\usepackage{subcaption}
\usepackage{float}
\usepackage{booktabs}
\graphicspath{{../figures/}}

% ---------------------------------------------------------------------------
% Colors
% ---------------------------------------------------------------------------
\usepackage{xcolor}

% ---------------------------------------------------------------------------
% Lists
% ---------------------------------------------------------------------------
\usepackage{enumitem}
\setlist[itemize]{nosep, topsep=0.3em, itemsep=0.2em}

% ---------------------------------------------------------------------------
% Section spacing
% ---------------------------------------------------------------------------
\usepackage{titlesec}
\titlespacing*{\section}{0pt}{1.2em}{0.6em}
\titlespacing*{\subsection}{0pt}{0.9em}{0.3em}

% ---------------------------------------------------------------------------
% Header and footer
% ---------------------------------------------------------------------------
\usepackage{fancyhdr}
\pagestyle{fancy}
\fancyhf{}
\lhead{\small Neural Surrogate for Snake Robot RL}
\rfoot{\thepage}
\renewcommand{\headrulewidth}{0.4pt}

% ---------------------------------------------------------------------------
% Table of contents depth
% ---------------------------------------------------------------------------
\setcounter{tocdepth}{2}
\setcounter{secnumdepth}{3}

% ---------------------------------------------------------------------------
% Placeholder macros
% ---------------------------------------------------------------------------
\newcommand{\placeholder}[1]{\textcolor{gray}{\textit{[Placeholder: #1]}}}
\newcommand{\mytodo}[1]{\textcolor{red}{\textbf{[TODO: #1]}}}

% ---------------------------------------------------------------------------
% Bibliography (natbib — before hyperref)
% ---------------------------------------------------------------------------
\usepackage[round, sort, authoryear]{natbib}

% ---------------------------------------------------------------------------
% Cross-references (cleveref — after amsmath, before hyperref)
% ---------------------------------------------------------------------------
\usepackage{cleveref}

% ---------------------------------------------------------------------------
% Hyperref and microtype — LAST
% ---------------------------------------------------------------------------
\usepackage[hidelinks]{hyperref}
\usepackage{microtype}

% ---------------------------------------------------------------------------
% Title
% ---------------------------------------------------------------------------
\title{%
  \textbf{Neural Surrogate Modeling for Snake Robot Locomotion:}\\[0.3em]
  \large Faster Reinforcement Learning via Learned Cosserat Rod Dynamics%
}
\author{[Author Name]}
\date{}

% ===========================================================================
\begin{document}
% ===========================================================================

\maketitle

% ---------------------------------------------------------------------------
\begin{abstract}
\placeholder{Write abstract after all results are in.}
\end{abstract}

\tableofcontents
\newpage

% ---------------------------------------------------------------------------
\section{Introduction}
\label{sec:intro}

\placeholder{Introduction: problem statement, contribution summary, paper roadmap.}

% ---------------------------------------------------------------------------
\section{Background}
\label{sec:background}

\subsection{Cosserat Rod Dynamics}
\label{sec:background:cosserat}

The body of the snake robot is modeled as a Cosserat rod --- a one-dimensional elastic continuum that supports shear, extension, bending, and torsion along its arc length \citep{Till2019}.
Unlike rigid-body chain models, the Cosserat formulation captures the distributed compliance of the rod, making it well-suited for soft and continuum robots whose deformation is governed by material elasticity rather than joint constraints.
The governing equations, stated per unit arc length $s$, are the linear and angular momentum balances:

\begin{align}
  \rho A \frac{\partial^2 \mathbf{x}}{\partial t^2}
    &= \frac{\partial \mathbf{F}}{\partial s} + \mathbf{f}_{\text{ext}},
  \label{eq:linear-momentum} \\[4pt]
  \rho \mathbf{I} \frac{\partial^2 \mathbf{d}}{\partial t^2}
    &= \frac{\partial \mathbf{M}}{\partial s} + \mathbf{m}_{\text{ext}} + \mathbf{x}' \times \mathbf{F},
  \label{eq:angular-momentum}
\end{align}

\noindent
where $\mathbf{x}(s,t) \in \mathbb{R}^3$ is the rod centerline position, $\mathbf{d}(s,t)$ is the material director frame encoding local orientation, $\mathbf{F}$ is the internal force (shear and tension, from the constitutive law), $\mathbf{M}$ is the internal moment (bending and torsion), $\mathbf{f}_{\text{ext}}$ denotes external contact forces from ground interaction, and $\mathbf{m}_{\text{ext}}$ represents the muscle moments generated by the CPG actuator.
The term $\rho A$ is the linear mass density and $\rho \mathbf{I}$ is the rotational inertia per unit length.

Although the full Cosserat model operates in three dimensions, snake locomotion in this work is constrained to the horizontal plane: all $z$-components of position and velocity are verified to be identically zero, so the system reduces to a two-dimensional projection without loss of physical content.

After spatial discretization into $N = 21$ nodes ($N_e = 20$ elements), the Cosserat PDEs (\cref{eq:linear-momentum,eq:angular-momentum}) become a first-order ordinary differential equation (ODE) system that is integrated forward in time by a symplectic Euler stepper.
The resulting discrete state vector $\mathbf{s}_t \in \mathbb{R}^{124}$ is laid out as follows:

\begin{table}[H]
  \centering
  \small
  \caption{Layout of the 124-dimensional rod state vector $\mathbf{s}_t$.}
  \label{tab:state-vector}
  \begin{tabular}{llrl}
    \toprule
    Index slice & Variable & Count & Physical meaning \\
    \midrule
    \texttt{[0:21]}   & $x_i$          & 21 & Node $x$-positions (m)                         \\
    \texttt{[21:42]}  & $y_i$          & 21 & Node $y$-positions (m)                         \\
    \texttt{[42:63]}  & $\dot{x}_i$    & 21 & Node $x$-velocities (m\,s$^{-1}$)             \\
    \texttt{[63:84]}  & $\dot{y}_i$    & 21 & Node $y$-velocities (m\,s$^{-1}$)             \\
    \texttt{[84:104]} & $\psi_e$       & 20 & Element yaw $= \arctan2(\text{tang}_y, \text{tang}_x)$ (rad) \\
    \texttt{[104:124]}& $\omega_{z,e}$ & 20 & Element angular velocity about $z$-axis (rad\,s$^{-1}$) \\
    \bottomrule
  \end{tabular}
\end{table}

\noindent
This 124-dimensional state is the Markov state used for both surrogate training and RL policy input.
The lower-dimensional observation space (14 floats based on FFT compression of node positions) that the policy receives is not Markov --- identical observations can yield different next-states depending on the hidden rod configuration and velocity history --- so the surrogate must operate on the full 124-dimensional rod state to produce correct multi-step rollouts.

\subsection{Resistive Force Theory Contact Model}
\label{sec:background:rft}

Ground contact forces are modeled using Resistive Force Theory (RFT), a slender-body approximation in which the contact force at each rod element depends only on the local velocity components resolved into directions tangent and normal to the element's centerline \citep{Till2019}.
This formulation avoids the computational cost of solving a full contact optimization at each substep, while capturing the anisotropic friction that is the physical mechanism of serpentine propulsion.

The tangential and normal drag forces per unit length acting on element $e$ are:
\begin{align}
  f_t &= -c_t \, v_t, \label{eq:rft-tangential}\\
  f_n &= -c_n \, v_n, \label{eq:rft-normal}
\end{align}

\noindent
where $v_t$ and $v_n$ are the velocity components tangent and normal to the element, respectively, and $c_t, c_n$ are anisotropic drag coefficients.
In this work we use $c_t = 0.01$ and $c_n = 0.05$, giving a normal-to-tangential ratio of five.
The anisotropy $c_n > c_t$ is the essential physical mechanism that enables net forward motion from a traveling curvature wave: lateral slip is resisted more strongly than axial slip, so the reaction forces from the ground produce a net thrust in the forward direction.

\subsection{Central Pattern Generator Actuation}
\label{sec:background:cpg}

The controller generates a traveling curvature wave $\kappa^{\text{target}}$ across the rod body; an elastic restoring force within PyElastica then drives each element toward its local target curvature, producing the characteristic serpentine gait.
This Central Pattern Generator (CPG) approach parameterizes the entire gait with a small number of physically meaningful variables, which facilitates both learning and analysis.

The target curvature at element $i$ is:
\begin{equation}
  \kappa_i^{\text{target}} = A \sin\!\bigl(2\pi k \cdot s_i - 2\pi f \cdot t + \phi_0\bigr) + b,
  \label{eq:cpg-curvature}
\end{equation}

\noindent
where $s_i \in [0, 1]$ is the normalized arc-length position of element $i$, $A$ is the amplitude (rad\,m$^{-1}$), $k$ is the wave number (wavelengths per body length), $f$ is the oscillation frequency (Hz), $\phi_0$ is the initial phase offset (rad), and $b$ is a constant turn bias (rad\,m$^{-1}$).
The five-dimensional action vector $\mathbf{a} = (A, f, k, \phi_0, b)$ parameterizes this wave, with physical ranges $A \in [0, 5]$, $f \in [0.5, 3.0]$\,Hz, $k \in [0.5, 3.5]$, $\phi_0 \in [0, 2\pi]$, and $b \in [-2, 2]$.
The RL policy outputs a normalized action $\mathbf{a} \in [-1,1]^5$ that is linearly mapped to these physical ranges at each control step.

\subsection{PyElastica Simulator}
\label{sec:background:elastica}

PyElastica \citep{Naughton2021} is an open-source Python framework for compliant mechanics simulation based on Cosserat rod theory.
It implements the discretized equations of motion described in \cref{sec:background:cosserat} using a symplectic Euler integrator, which provides numerical stability for the stiff elastic rod dynamics without the computational cost of implicit methods.
Each RL control step (of duration $\Delta t_{\text{ctrl}}$) runs $n_{\text{sub}}$ internal physics substeps, defining the ground-truth transition operator $T : (\mathbf{s}_t, \mathbf{a}_t) \mapsto \mathbf{s}_{t+1}$.

Each evaluation of $T$ requires approximately 46\,ms when running 16 parallel environment workers on a 48-core CPU server, yielding a throughput of roughly 350 environment steps per second.
At this rate, a PPO training run requiring $2 \times 10^6$ environment interactions takes approximately 36 hours of wall-clock time.
This computational bottleneck motivates replacing $T$ with a neural surrogate model $f_\theta \approx T$ that can be evaluated in a single GPU forward pass (${\sim}0.001$\,ms per step), enabling massive parallelization and a projected throughput improvement of four to five orders of magnitude.

% ---------------------------------------------------------------------------
\section{Related Work}
\label{sec:related}

\subsection{Neural Surrogates for Soft Robots}
\label{sec:related:surrogates}

Approximating the dynamics of elastic and soft robots with neural networks has attracted substantial research attention as a means of bypassing the prohibitive cost of first-principles simulation.
The strongest result in this direction for Cosserat rod dynamics comes from \citet{Stolzle2025}, who propose a domain-decoupled physics-informed neural network (DD-PINN) that achieves a 44,000$\times$ speedup over numerical Cosserat rod integration while maintaining end-effector position errors below 3\,mm (2.3\% of actuator length) across the full 72-dimensional rod state.
The key innovation is the decoupling of the time domain from the feed-forward network, which enables gradients to be computed in closed form rather than via expensive automatic differentiation through the ODE integrator.
The resulting model enables nonlinear model-predictive control (MPC) at 70\,Hz on a GPU --- a capability that is simply infeasible with direct numerical integration.

Building on this foundation, \citet{Hsieh2024} adopt a hybrid approach in KNODE-Cosserat: a simplified first-principles Cosserat rod model serves as the backbone, while a neural ODE residual term corrects the systematic error between the simplified physics and real hardware behavior.
Crucially, the training data for KNODE-Cosserat is collected from PyElastica simulations --- a pipeline directly analogous to the one used in this work.
The resulting hybrid model achieves a 58.7\% accuracy improvement over the physics-only model on long-horizon trajectories, demonstrating that even a coarse physics prior substantially reduces the burden on the neural correction term.

The most directly relevant prior work to this paper is SoRoLEX \citep{SoRoLEX2024}, which trains an LSTM forward dynamics model on soft robot simulator data and uses the learned model as a surrogate environment for PPO training with JAX-based GPU parallelization.
The central pipeline --- collect trajectories from a high-fidelity simulator, train a neural dynamics model, and run large-scale RL in the learned environment --- is identical to the approach taken here.
Where this work differs is in the choice of architecture (MLP rather than LSTM) and in the prediction target (the full 124-dimensional Cosserat rod state rather than a task-space observation), and in the explicit use of a multi-step rollout loss to limit autoregressive drift during RL rollouts.

In a related thread, \citet{PINNSoftRobot2025} demonstrate that physics-informed neural networks for articulated soft robots can achieve a 467$\times$ speedup over accurate first-principles models while enabling MPC at 47\,Hz, and require as little as a single training dataset --- illustrating the data efficiency advantages that come from incorporating physical structure into the surrogate architecture.

Taken together, these results establish that neural surrogates for Cosserat rod and soft robot dynamics are feasible, practical, and capable of sustaining the multi-step rollout accuracy required for closed-loop control.
This work follows the data-driven approach of SoRoLEX but with an MLP predicting the full Cosserat rod state (124 dimensions), trained with both single-step MSE and multi-step rollout loss to reduce the autoregressive error accumulation that poses the main challenge for long-horizon RL training.

\subsection{Reinforcement Learning with Learned Simulators}
\label{sec:related:rl}

The idea of replacing expensive environment interactions with a learned dynamics model dates to Dyna \citep[Sutton, 1990, as cited in][]{Janner2019}, and has recently been placed on a rigorous theoretical and empirical footing by a series of influential results.
\citet{Janner2019} show in Model-Based Policy Optimization (MBPO) that an ensemble of learned dynamics models can be used to generate short synthetic rollouts from real environment states, and that policy optimization on these rollouts achieves model-free asymptotic performance with approximately five times less real-environment interaction.
The central insight of MBPO --- that short-horizon model predictions limit error accumulation while still providing useful policy gradient signal --- directly motivates the use of a surrogate environment for RL in this work.

The world model paradigm extends this approach by learning compact latent representations of environment dynamics alongside the policy.
\citet[DreamerV3,][]{Janner2019} demonstrate that a recurrent state-space model trained entirely from sensory observations can sustain imagination-based RL across 150 diverse tasks without task-specific tuning, illustrating that surrogate quality is the primary bottleneck for downstream policy performance.
In the context of the present work, the accuracy of the Cosserat rod surrogate --- quantified in Phase 4 validation --- directly determines the quality of policies trained in the surrogate environment.
Our Phase 8 comparison with a direct PyElastica RL baseline isolates the policy gap attributable to surrogate modeling error.

Graph-structured alternatives to MLP surrogates have also been explored for deformable body simulation.
\citet{SanchezGonzalez2020} introduce a particle-based graph network simulator (GNS) that treats simulation as message passing over a particle graph, achieving up to 1000$\times$ speedups for fluid, rigid, and deformable material simulation.
\citet{Pfaff2021} extend this approach to mesh-based simulation with MeshGraphNets, which operate directly on unstructured simulation meshes and generalize across different mesh resolutions.
For the snake robot considered in this work, however, the rod is discretized into only 20 elements --- a sufficiently small graph that the added architectural complexity of a GNN is not warranted, and a dense MLP achieves comparable expressive power at a fraction of the implementation cost.
Graph-structured surrogates would become advantageous if the problem were extended to rods with significantly more nodes or to multi-rod assemblies interacting through contact.

% ---------------------------------------------------------------------------
\section{Methods}
\label{sec:methods}

\subsection{Data Collection Pipeline}
\label{sec:methods:datacollection}

\placeholder{Data collection pipeline: parallel PyElastica workers, Sobol quasi-random action sampling, 30\% perturbation, 50\% random action fraction, checkpoint-format V2 data (5 boundary states per run), 25 GB target.}

\subsection{State and Action Representation}
\label{sec:methods:representation}

\subsection{Per-Element CPG Phase Encoding}
\label{sec:methods:encoding}

\placeholder{Per-element CPG phase encoding: why raw serpenoid\_time is non-Markov (unbounded), sin/cos encoding $(\sin(\omega t_i), \cos(\omega t_i))$ per element, $\phi_i$ computed on-the-fly from action parameters and $t_{\text{start}}$, INPUT\_DIM expanded from 131 to 189.}

\subsection{Surrogate Model Architecture}
\label{sec:methods:architecture}

\placeholder{Surrogate architecture: MLP with state-delta prediction, input $\in \mathbb{R}^{189}$ = rod state (124) + action (5) + per-element phase encoding (60), output $\in \mathbb{R}^{124}$ (state delta), hidden layers $512 \times 3$, SiLU activations, residual connection $s_{t+1} = s_t + \Delta$.}

\subsection{Training Procedure}
\label{sec:methods:training}

\placeholder{Training procedure: MSE loss on state deltas, multi-step rollout loss, learning rate schedule, batch size, number of epochs, W\&B logging, architecture sweep configurations.}

% ---------------------------------------------------------------------------
\section{Experiments and Results}
\label{sec:experiments}

\subsection{Surrogate Accuracy}
\label{sec:experiments:accuracy}

\placeholder{Phase 4/5/8 results here. Surrogate validation against held-out PyElastica trajectories: 1-step MSE, multi-step rollout RMSE, per-state-component accuracy ($\omega_z$ poor prediction issue and fix).}

\subsection{RL Training with Surrogate}
\label{sec:experiments:rl}

\placeholder{Phase 4/5/8 results here. PPO training on surrogate environment: learning curves, wall-clock time comparison, policy performance on surrogate vs real Elastica environment.}

\subsection{Comparison with Direct Elastica RL Baseline}
\label{sec:experiments:baseline}

\placeholder{Phase 4/5/8 results here. Head-to-head comparison: surrogate RL vs direct Elastica RL (Phase 8 baseline). Sample efficiency, wall-clock time, final policy quality, transfer gap.}

% ---------------------------------------------------------------------------
\section{Discussion}
\label{sec:discussion}

\subsection{Physics Calibration Challenges}
\label{sec:discussion:physics}

\placeholder{Physics calibration challenges: rod radius (0.001 m $\rightarrow$ 0.02 m, $EI \propto r^4$), Young's modulus tuning (2e6 $\rightarrow$ 1e5 Pa), RFT friction coefficients ($c_t$, $c_n$), serpenoid wave direction sign error, damping parameter. These were critical blockers for stable locomotion.}

\subsection{Surrogate Architecture Experiments}
\label{sec:discussion:arch}

\placeholder{Surrogate architecture experiments: rollout loss variants, residual MLP vs baseline MLP, history window experiments, architecture sweep comparison findings (Phase 03 Plans 03--05).}

\subsection{Data Collection Pipeline Challenges}
\label{sec:discussion:pipeline}

\placeholder{Data collection pipeline challenges: parallel scaling bottleneck (IPC overhead dominates at $>$16 workers), Numba thread pool deadlock, stall detection false positives during init, omega\_z coverage gap (perturb\_omega\_std: 0.05 $\rightarrow$ 1.5 rad/s fix), data loss on worker respawn.}

% ---------------------------------------------------------------------------
\section{Conclusion}
\label{sec:conclusion}

\placeholder{Write conclusion after all results are in.}

% ---------------------------------------------------------------------------
% Bibliography
% ---------------------------------------------------------------------------
\bibliographystyle{plainnat}
\bibliography{references}

\end{document}
